\section{Diseño de la PCB}

\subsection{Apilado y contorno}

La placa tiene \bLayers{} capas de cobre sobre FR4 de \bThickness~mm (el espesor del diseño
recomendado del conector USB-C de 16 pines), con un contorno de \bBoardWidth{} $\times$
\bBoardHeight~mm. La capa inferior lleva un vertido de masa continuo, que sirve de plano de
referencia de las líneas de radiofrecuencia y del USB. Al ser de dos capas, el fabricante no ofrece
impedancia controlada, y el diseño acepta esa dispersión en líneas cortas.

\subsection{Clases de red y reglas cargadas en Altium}

\begin{table}[H]
\centering
\footnotesize
\begin{tabularx}{\textwidth}{|l l X|}
\hline
\tablehead{\textbf{Clase} & \textbf{Ancho preferido} & \textbf{Uso}} \\ \hline
\texttt{RF50} & \rRfWidth~mm & Línea de 50~$\Omega$ del LoRa y del GNSS hasta sus conectores u.FL, microstrip sobre el plano de la capa inferior \\ \hline
\texttt{USB} & \rUsbWidth~mm & Par diferencial D+/D$-$, con longitudes igualadas \\ \hline
\ruta{PWR_CN3791} & \rChargerWidth~mm & Lazo del cargador solar (diodos, MOSFET, bobina) \\ \hline
\texttt{PWR} & \rPwrWidth~mm & 3V3, batería y alimentación en general \\ \hline
Resto & 0.203~mm & Señales \\ \hline
\end{tabularx}
\caption{Clases de red y ancho de pista preferido.}
\end{table}

Reglas generales: separación mínima entre cobre de \rClearance~mm, vías con agujero de \rViaHole~mm y
diámetro de \rViaDiameter~mm, y separaciones propias del componente USB-C y del U8. El ancho de la línea de
50~$\Omega$ salió de un cálculo de microstrip sobre FR4 de \bThickness~mm con $\varepsilon_r = 4.4$
(estimación, \ruta{hardware/PCB_REGLAS.md}): con la tolerancia de espesor del fabricante queda entre
unos 46 y 54~$\Omega$.

\subsection{Criterios de ubicación}

Los criterios de zonificación salen de las guías de integración de u-blox y de Espressif
(\ruta{hardware/PCB_REGLAS.md}):
\begin{itemize}
  \item los componentes de radiofrecuencia del GNSS, a 5~mm como mínimo de cualquier otro componente de RF y lejos de lo digital;
  \item el parche GNSS y su conector, del lado opuesto al Wi-Fi del módulo y al LoRa;
  \item el MAX-M10S, lejos de fuentes de calor (TP4056, CN3791, bobina) y sin pistas de otras señales debajo del módulo;
  \item el lazo del CN3791 lo más corto posible, y su etapa de potencia lo más lejos posible del GNSS y del cristal;
  \item el pad térmico del TP4056 a masa con una zona amplia de cobre y vías hacia la capa inferior.
\end{itemize}
Que el layout final cumpla cada criterio no se ha auditado para esta nota; es un punto de la revisión
de la placa antes de fabricarla (sección~\ref{sec:limitaciones}).

\subsection{Verificación de reglas de diseño}

El DRC de Altium del 15 de noviembre de 2026 comprobó \bDrcRules{} reglas y detectó \bDrcViolations{}
violaciones, sin ninguna excepción concedida. Entre ellas: separación de cobre, cortocircuitos, nets sin
rutear, anchos por clase, vías, anillo mínimo, separación entre agujeros, lengüetas de máscara, serigrafía,
longitudes igualadas del par USB y separaciones entre componentes. El resumen de cada regla está en
\ruta{evidence/drc_summary.txt}. El DRC verifica las reglas que se cargaron, no el diseño completo: no
sustituye a la revisión de integridad de señal ni a la de radiofrecuencia.

\subsection{Taladros}

\begin{table}[H]
\centering
\small
\begin{tabular}{|l l r l|}
\hline
\tablehead{\textbf{Herramienta} & \textbf{Diámetro} & \textbf{Cantidad} & \textbf{Tipo}} \\ \hline
T1 & 0.305~mm & 160 & Pasante metalizado \\ \hline
T2 & 0.381~mm & 11 & Pasante metalizado \\ \hline
T3 & 0.6~mm & 4 & Ranura metalizada \\ \hline
T4 & 0.762~mm & 4 & Pasante metalizado \\ \hline
T5 & 0.8~mm & 2 & Pasante metalizado \\ \hline
T6 & 0.85~mm & 2 & Pasante metalizado \\ \hline
T7 & 0.65~mm & 2 & No metalizado \\ \hline
\end{tabular}
\caption{Taladros de la placa (informe de taladro de Altium); \bHoles{} en total.}
\end{table}
